\chapter{Metodología y planificación}
\label{chap:metodologia}

\lettrine{E}{n} este capítulo se describe la metodología seguida durante la
realización del proyecto y su planificación.
Además, se detallan los conjuntos de datos empleados en el
entrenamiento y la evaluación de los modelos.

\section{Enfoque Metodológico}
Al tratarse de un trabajo de investigación con un fuerte componente
experimental, se adoptó un enfoque iterativo e incremental. Esta metodología
permitió construir el sistema de forma gradual, validando cada técnica antes de
incorporar la siguiente y añadiendo valor en cada iteración. A lo largo del
proceso se siguieron los principios de la metodología ágil \textit{Scrum}~\cite{schwaber2020scrum},
basada en ciclos de trabajo cortos y repetitivos
({\glspl{sprint}}) de una duración media de una semana.

Al final de cada iteración se realizaba una evaluación del progreso,
identificando los problemas surgidos y proponiendo soluciones para
continuar con el desarrollo. Esta práctica resultó especialmente útil en un
proyecto con un alto grado de incertidumbre experimental, en el que el resultado de
una técnica condicionaba el diseño de la siguiente. Cada \gls{sprint}
concluía con una reunión con los directores del proyecto, en la
que se revisaban los avances y se establecían los siguientes pasos.

\section{Fases del Proyecto}
El proyecto se inspira fuertemente en el estado del arte (Capítulo~\ref{chap:sota}), con el
que comparte algunos conceptos de diseño (conjunto de acciones, estados, recompensas\dots), pero del que 
difiere en su implementación al probar familias de técnicas muy variadas entre sí, incluyendo algoritmos
no implementados directamente en \textit{Minecraft}, como las \gls{htn}. A continuación se presentan tanto la planificación
inicial como las fases finales en las que se dividió el trabajo.

\subsection{Fases \textit{a priori}}
Antes de comenzar el desarrollo se definieron las siguientes cuatro fases:

\begin{enumerate}
  \item \textbf{Análisis y preparación inicial}: estudio del estado del arte
        sobre agentes en \textit{Minecraft} y familiarización con las herramientas
        (entorno de simulación, comunicación con el agente y
        {\glspl{framework}}).
  \item \textbf{Planificación simbólica}: implementación de la
        \gls{htn} como primera técnica y \textit{baseline} interpretable
        para la comparativa.
  \item \textbf{Aprendizaje por refuerzo}: implementación de un agente de
        refuerzo capaz de resolver la tarea únicamente a partir de la
        observación del entorno.
  \item \textbf{Ejecución y evaluación final}: ejecución de las técnicas,
        configuración necesaria y análisis comparativo de los resultados.
\end{enumerate}

\subsection{Fases \textit{a posteriori}}
En el desarrollo real, el trabajo se dividió en siete iteraciones. Debido a que la aproximación
incremental, el diseño y la realización de cada fase condicionaban la siguiente. Por ejemplo,
la fase de \gls{htn} se usó para generar un conjunto de demostraciones 
que sirvió como base para el \gls{il}, el cual se añadió con posterioridad a la planificación inicial.

\begin{enumerate}
  \item \textbf{Análisis y preparación inicial}: revisión del trabajo previo,
        junto con la familiarización con el
        entorno de simulación y el diseño de la arquitectura general del
        sistema, descrito en el Capítulo~\ref{chap:arquitectura}.
  \item \textbf{Entorno de simulación}: implementación de un entorno
        \gls{gymnasium} comunicado mediante \acrshort{http} con un agente \textit{Mineflayer}~\cite{mineflayer}
        sobre un servidor real de \textit{Minecraft}, con reinicio determinista.
  \item \textbf{Planificación simbólica y generación del \textit{dataset}}:
        implementación del planificador \gls{htn} y de la descomposición de la tarea
        de \gls{woodcutting}, empleando sus ejecuciones como conjunto de
        demostraciones para la fase de \gls{il}.
  \item \textbf{Aprendizaje por imitación}: diseño iterativo y entrenamiento final de una \gls{política} con
        arquitectura de doble cabeza sobre el \textit{dataset} generado por el
        \gls{htn}.
  \item \textbf{Aprendizaje por refuerzo}: diseño iterativo y entrenamiento de una \gls{política} de refuerzo,
        optimizando la arquitectura y los hiperparámetros.
  \item \textbf{Evaluación}: introducción de métricas, definición de una condición de éxito explícita y
        ajuste de la función de recompensa.
  \item \textbf{Resultados y evaluación final}: evaluación de los resultados de todas las implementaciones,
        analizando sus fortalezas y debilidades.
\end{enumerate}

Durante el desarrollo se exploraron además otras dos aproximaciones que
finalmente se descartaron del alcance de la comparativa. La primera fue el uso de
agentes basados en \glspl{llm} mediante el \gls{framework}
\textit{Mindcraft}~\cite{mindcraft2025}, descartada por el elevado coste computacional de ejecutar \glspl{llm} pesados en local. 
La segunda fue un detector visual de árboles
basado en \acrshort{yolo}~\cite{DBLP:conf/cvpr/RedmonDGF16}, descartado por desplazar el objetivo 
hacia un problema de detección de objetos en lugar del aprendizaje de la tarea en sí.

\section{Recursos y limitaciones}
Para el desarrollo del proyecto se emplearon dos equipos: un ordenador de
sobremesa, encargado de ejecutar los entrenamientos, alojar el
servidor de \textit{Minecraft}, ejecutar el agente y el desarrollo general; y un portátil MacBook Air M3,
empleado como equipo secundario. Sus características se resumen en la Tabla~\ref{tab:hardware}
(página~\pageref{tab:hardware}).

\begin{table}[htbp]
  \centering
  \resizebox{\textwidth}{!}{
  \small
  \rowcolors{2}{white}{udcgray!25}
  \begin{tabular}{c|cc}
  \rowcolor{udcpink!25}
  \textbf{Característica} & \textbf{Sobremesa} & \textbf{MacBook Air M3} \\\hline
  \acrshort{cpu} & AMD Ryzen 9 5900X (\num{12}\,C/\num{24}\,T) & Apple M3 (\num{8} núcleos) \\
  \acrshort{ram} & \qty{32}{GB} & \qty{16}{GB} \\
  \acrshort{gpu} & NVIDIA RTX 3060 \qty{12}{GB} & Integrada (\num{10} núcleos) \\
  Uso & Entrenamiento y desarrollo & Equipo complementario \\
  \end{tabular}
  }
  \caption{Equipos empleados en el desarrollo del proyecto.}
  \label{tab:hardware}
\end{table}

La principal limitación de recursos fue el tiempo de entrenamiento. Si bien el
sistema no generaba una gran carga computacional, la ejecución de las acciones
suponía una limitación temporal. El entorno de simulación era el cuello de
botella a la hora de procesar las acciones. Si el agente realiza la acción de
talar y esta requiere \numrange{2}{3} segundos para completarse, en un episodio de \num{100}
pasos el tiempo de ejecución asciende a \numrange{200}{300} segundos. Considerando la ejecución
completa, con \num{300} episodios se alcanzan entre \num{16} y \num{25} horas de entrenamiento, lo
que limita el número de experimentos que se pueden realizar.

La única paralelización posible sería ejecutar varios agentes a la vez o usar
varios equipos simultáneamente, pero no se hizo debido a la limitación de recursos
computacionales disponibles y a la dificultad de coordinar el desarrollo en
varios equipos a la vez.

\section{Conjuntos de Datos}
A diferencia de un problema de clasificación supervisada con conjuntos de datos
estáticos, en este trabajo los datos provienen de la interacción de agentes con
el entorno. Se consideraron dos conjuntos de demostraciones, seleccionando el
primero como conjunto final.

\paragraph{\textit{Dataset} \gls{htn}:} este conjunto de datos de elaboración propia se compone de las 
trayectorias producidas por el agente simbólico al resolver la tarea de
\textit{woodcutting}. Estas trayectorias se registran como pares \textit{OBSERVACIÓN-ACCIÓN} 
(definidos en el Capítulo~\ref{chap:arquitectura}) y conforman el
conjunto de demostraciones de un experto determinista e interpretable. Este
\textit{dataset} permite la comparación entre el comportamiento simbólico y las
técnicas de aprendizaje profundo. Se realiza una pequeña limpieza descartando las 
muestras de acciones poco representativas, quedando un total de \num{118} episodios y
\num{4094} \textit{steps} (\num{4069} tras la limpieza), como se detalla en la Tabla~\ref{tab:datasets} (página~\pageref{tab:datasets}).

\paragraph{\textit{Dataset} MineRL Treechop-v0:} este conjunto de demostraciones humanas~\cite{guss2019minerl}
se compone de \num{210} episodios (todos con éxito) y aproximadamente \num{453000}
\textit{steps}, con una mediana de \num{64} troncos recogidos por episodio.
Su análisis exploratorio reveló que el movimiento de cámara humano (mediana de
\ang{0.75}/\ang{0.90} por \textit{step}) es aproximadamente diez veces
más fino que el paso de cámara de la discretización empleada por los agentes discretos, lo que impide trasladar
directamente las demostraciones humanas a su espacio de acción. Además, el \textit{dataset} humano
presenta limitaciones al ser utilizado fuera del entorno de desarrollo de \textit{MineRL}.
Entre estas limitaciones se encuentran: la estructura de los datos, el formato de las acciones y 
el formato de las observaciones, lo que dificulta su uso directo para el
entrenamiento de agentes en nuestro entorno de simulación. Por último, \textit{MineRL} solo proporciona información
visual, por lo que no permite obtener el vector de estado ni las señales de percepción
(\textit{tree\_visible}, \textit{tree\_distance}) de las que dependen nuestros agentes. 

\begin{table}[htbp]
  \centering
  \resizebox{\textwidth}{!}{
  \small
  \rowcolors{2}{white}{udcgray!25}
  \begin{tabular}{c|cc}
  \rowcolor{udcpink!25}
  \textbf{Propiedad} & \textbf{{htn}} & \textbf{MineRL Treechop-v0} \\\hline
  Origen         & Planificador simbólico (propio) & Demostraciones humanas \\
  Episodios      & \num{118} & \num{210} \\
  \textit{Steps} (bruto)  & \num{4094} & $\sim$\num{453000} \\
  \textit{Steps} (limpio) & \num{4069} & $\sim$\num{453000} \\
  \end{tabular}
  }
  \caption[Conjuntos de datos analizados]{Comparación de los conjuntos de datos analizados en el trabajo.}
  \label{tab:datasets}
\end{table}

Cabe mencionar además que el \gls{framework} de \textit{MineRL} se intentó utilizar, pero actualmente ya no recibe mantenimiento y no es compatible con las versiones actuales de \gls{gymnasium}~\cite{minerl_issue788}. Por otra parte,
este tipo de \textit{datasets} están pensados para ser utilizados dentro de su propio \gls{framework}, 
lo que dificulta su uso en entornos de simulación externos como el empleado en este proyecto.